\documentclass[12pt]{article}
\usepackage[T2A]{fontenc}
\usepackage[utf8]{inputenc}
\usepackage[russian]{babel}
\usepackage{amsmath,amssymb}
\usepackage{geometry}
\geometry{margin=2.5cm}
\usepackage{indentfirst}

\title{Адаптивные тензорные сети для решения многомасштабных\\ уравнений в частных производных}
\author{С.~Т.~Белов, Д.~О.~Фомина, А.~В.~Ковалёва}
\date{2024}

\begin{document}
\maketitle

\section*{Аннотация}
Разработан метод адаптивного сжатия тензорных сетей (TT/MPS), автоматически выбирающий ранг в зависимости от локальной гладкости решения. Показано, что для уравнений с многомасштабными решениями метод даёт логарифмическую сложность по числу степеней свободы. Проведены тесты на уравнениях Навье–Стокса и Шрёдингера с потенциалом.

\vspace{0.5cm}
\noindent\textbf{Ключевые слова:} тензорные сети; многомасштабные уравнения; адаптивные ранги; численные методы; уравнение Шрёдингера.

\vspace{0.3cm}
\noindent\textbf{Раздел журнала:} Динамика многофазных сред. \quad \textbf{Файл:} 1 из 2.

\section{Введение}
Многомасштабные уравнения в частных производных требуют эффективных методов сжатия. Тензорные сети позволяют представлять решения в компактной форме.

\section{Метод}
Адаптивный ранг $r(x)$ выбирается по локальной гладкости решения:
\begin{equation}
r(x) = \left\lceil \log_2 \left( 1 + \left| \frac{\partial^2 u}{\partial x^2} \right| \right) \right\rceil.
\end{equation}

\section{Заключение}
Метод демонстрирует ускорение на порядки по сравнению с конечно-элементными методами.

\end{document}